% Generated by scripts/assemble.py; corrections belong in source/corrections.json.
\examhead{97}{南京理工大学 2026 年数学分析试卷(回忆版)}{96}
% p097-b17; source PDF page 97
\blockanchor{p097-b17}
(注.题目来自考生回忆,部分表述或数据可能有误)\par
% p097-b18; source PDF page 97
\blockanchor{p097-b18}
\question{1}解答如下问题:\par
% p097-b19; source PDF page 97
\blockanchor{p097-b19}
(1) 利用定义证明:  $(\cos x)' = -\sin x$ .\par
% p097-b20; source PDF page 97
\blockanchor{p097-b20}
(2) 利用定积分定义证明: $\int_{0}^{1} x^{3} \, \mathrm{d}x = \frac{1}{4}$.\par
% p097-b21; source PDF page 97
\blockanchor{p097-b21}
\question{2}求极限:\par
% p097-b22; source PDF page 97
\blockanchor{p097-b22}
(1) $\lim_{n\to \infty}\sum_{k = 1}^{n}\frac{2^{\frac{k}{n}}}{n + \frac{1}{k}}.$\par
% p097-b23; source PDF page 97
\blockanchor{p097-b23}
(2) (可能有误) $\lim_{x\to 0}\left(\frac{1}{\sin x} -\frac{\sin x}{e^x - 1}\right).$\par
% p097-b24; source PDF page 97
\blockanchor{p097-b24}
\question{3}设 $f(x, y) = \begin{cases} \frac{xy^2}{x^2 + y^4}, & (x, y) \neq (0, 0), \\ 0, & (x, y) = (0, 0). \end{cases}$ 讨论 $f(x, y)$ 在点 $(0, 0)$ 的连续性, 偏导数的存在性, 可微性.\par
% p097-b25; source PDF page 97
\blockanchor{p097-b25}
\question{4}求幂级数 $\sum_{n=0}^{\infty} \frac{(n+1)(n+2)}{n+3} x^n$ 的收敛域以及和函数.\par
% p098-b00; source PDF page 98
\blockanchor{p098-b00}
\question{5}设 $f(x)$ 在 $I$ 上有连续的导数, 记 $F_{n}(x) = nf\left(x + \frac{1}{n}\right) - nf(x)$.\par
% p098-b01; source PDF page 98
\blockanchor{p098-b01}
(1) 若 $I$ 为有界闭区间, $\{F_n(x)\}$ 是否一致收敛? 给出理由.\par
% p098-b02; source PDF page 98
\blockanchor{p098-b02}
(2) 若 $I$ 为有界开区间, $\{F_n(x)\}$ 是否一致收敛? 给出理由.\par
% p098-b03; source PDF page 98
\par\Needspace{6\baselineskip}
\blockanchor{p098-b03}
\question{6}设 $f(x)$ 为连续函数, 求曲线积分\par
\nopagebreak[4]
% p098-b04; source PDF page 98
\blockanchor{p098-b04}
\displayformula{\int_ {L} \frac {1 + y ^ {2} f (x y)}{y} \mathrm{d} x + \frac {x (y ^ {2} f (x y) - 1)}{y ^ {2}} \mathrm{d} y.}
% p098-b05; source PDF page 98
\blockanchor{p098-b05}
其中 L 为  $\left(3,\frac{2}{3}\right)$  指向 (1,2) 的直线.\par
% p098-b06; source PDF page 98
\blockanchor{p098-b06}
\question{7}证明: $\int_{0}^{+\infty} \frac{\sin x}{x} \, \mathrm{d}x = \frac{\pi}{2}$, 并求 $\int_{0}^{+\infty} \frac{\sin^2(xy)}{x^2} \, \mathrm{d}x$.\par
% p098-b07; source PDF page 98
\par\Needspace{6\baselineskip}
\blockanchor{p098-b07}
\question{8}求曲面积分\par
\nopagebreak[4]
% p098-b08; source PDF page 98
\blockanchor{p098-b08}
\displayformula{\oiint_ {\Omega} \frac {x}{(x ^ {2} + y ^ {2} + z ^ {2}) ^ {\frac {3}{2}}} \mathrm{d} y \mathrm{d} z + \frac {y}{(x ^ {2} + y ^ {2} + z ^ {2}) ^ {\frac {3}{2}}} \mathrm{d} z \mathrm{d} x + \frac {z}{(x ^ {2} + y ^ {2} + z ^ {2}) ^ {\frac {3}{2}}} \mathrm{d} x \mathrm{d} y.}
% p098-b09; source PDF page 98
\blockanchor{p098-b09}
其中 $\Omega$ 是区域 $|x| \leq 3, |y| \leq 5, |z| \leq 4$ 的表面, 方向向外.\par
% p098-b10; source PDF page 98
\blockanchor{p098-b10}
\question{9}解答如下问题:\par
% p098-b11; source PDF page 98
\blockanchor{p098-b11}
(1) 叙述有限覆盖定理以及一致连续的定义.\par
% p098-b12; source PDF page 98
\blockanchor{p098-b12}
(2) 利用有限覆盖定理证明连续函数在闭区间内一致连续.\par
